% 张三慧《大学物理学（第4版）》第一卷，第11章《热力学第二定律》。
% 源PDF物理页398--400，章末习题11.1--11.13，与思考题11.1--11.10严格区分。
% 此为源书事实重述的学术候选录入；图11.14、11.16后续须逐图重绘审核。
\section*{第11章\quad 热力学第二定律}
\subsection*{章末习题11.1--11.13}

\paragraph{习题 11.1}
一摩尔\emph{氧气}按\emph{刚性双原子理想气体}
（平动及转动5个自由度）处理，沿图11.14的
\(p\)-\(V\) 折线路径 \(a\to b\to c\) 演化。
图中纵轴单位为 \(10^5\,\mathrm{Pa}\)，横轴单位 \(\mathrm L\)；
按图点 \(a=(1\,\mathrm L,6\times10^5\,\mathrm{Pa})\)、
\(b=(2\,\mathrm L,8\times10^5\,\mathrm{Pa})\)、
\(c=(3\,\mathrm L,4\times10^5\,\mathrm{Pa})\)。
求全程气体对外做功、吸收的热量及熵变。
\textit{两段均按图中直线相连、准静态过程理解。}

\paragraph{习题 11.2}
将 \(30\,\mathrm g\)、\(-40^\circ\mathrm C\) 的冰，
经加热变成 \(100^\circ\mathrm C\) 的水蒸气。
冰和液态水比热分别取
\(c_i=2.1\,\mathrm{J/(g\,K)}\)、
\(c_w=4.2\,\mathrm{J/(g\,K)}\)，
水熔化潜热
\(\lambda=334\,\mathrm{J/g}\)，
汽化潜热 \(L=2260\,\mathrm{J/g}\)；
过程在 \(1.013\times10^5\,\mathrm{Pa}\) 下进行。
求水系统所需热量与熵变。

\paragraph{习题 11.3}
一个人一天约向周围环境散失热量
\(8\times10^5\,\mathrm J\)，
若周围环境按 \(273\,\mathrm K\) 的恒温热库处理，
估算该过程环境每天增加多少熵。
\textit{原题讨论的是向环境放出的热所引起的熵变。}

\paragraph{习题 11.4}
冬季房屋向外散热的功率
\(2\times10^8\,\mathrm{J/h}\)，
室内 \(20^\circ\mathrm C\)，
室外 \(-20^\circ\mathrm C\)。
求经这种温差传热产生的熵的速率，
单位为 \(\mathrm{J/(K\,s)}\)。

\paragraph{习题 11.5}
汽车匀速行驶时在摩擦中耗散的功率
\(20\,\mathrm{kW}\)。
气温为 \(12^\circ\mathrm C\)，
估计汽车摩擦所产生的熵的速率
\(\mathrm{J/(K\,s)}\)。

\paragraph{习题 11.6}
贵州黄果树瀑布主瀑布宽 \(83\,\mathrm m\)，
落差 \(74\,\mathrm m\)，
洪水期间流量
\(1500\,\mathrm{m^3/s}\)。
环境温度 \(12^\circ\mathrm C\)，
估算瀑布每秒由于水的重力势能
耗散为热而增加的熵。
水的密度近似 \(1000\,\mathrm{kg/m^3}\)。

\paragraph{习题 11.7}
（1）将 \(1\,\mathrm{kg}\)、\(0^\circ\mathrm C\)
的水放到 \(100^\circ\mathrm C\) 的恒温热库中，
直到水达到热平衡，计算“水+热库”系统的熵变，
并判定符号。
（2）改为先与 \(50^\circ\mathrm C\)
热库平衡，再与 \(100^\circ\mathrm C\)
热库平衡；重新计算整个系统的熵变
并与（1）比较。
取液态水比热 \(c=4.2\times10^3\,\mathrm{J/(kg\,K)}\)。

\paragraph{习题 11.8}
绝热良好的金属筒内起初
有 \(2.5\,\mathrm{kg}\) 水和
\(0.7\,\mathrm{kg}\) 冰，在 \(0^\circ\mathrm C\)
下处于冰水共存平衡。
（1）把筒置于温度\emph{略低于}
\(0^\circ\mathrm C\) 的房间，
使筒内冰与水的\emph{质量变成相等}，
求此过程混合物的熵变及
它与房间构成的系统的总熵变。
（2）再将筒置于 \(100^\circ\mathrm C\)
的恒温热库，令冰水混合物
\emph{恢复起初2.5kg水、0.7kg冰的状态}，
求这一过程冰水混合物与
混合物—高温热库系统各自熵变。
\textit{该题相变发生于接近0℃，
冻结或熔解的质量为0.9kg，不是全部相变。}

\paragraph{习题 11.9}
一定量理想气体初态
\(T_1=300\,\mathrm K\)，
\(p_1=3.039\times10^5\,\mathrm{Pa}\)，
\(V_1=4\,\mathrm{m^3}\)。
先等温膨胀到 \(V_2=16\,\mathrm{m^3}\)，
继而等容改变压强，
最后绝热压缩回初态。
所有过程可逆，比热比 \(\gamma=1.4\)。
（1）画出 \(p\)-\(V\) 循环及方向；
（2）求各段热量、对外功和熵变，
并求整个循环的数值。

\paragraph{习题 11.10}
绝热容器内将两份\emph{相同液体}
等质量 \(m\) 混合，初温分别为
\(T_1,T_2>T_1\)，
其定压比热 \(c_p\) 视为常数。
混合后达到共同终温，
求系统总熵变并证明严格大于零。

\paragraph{习题 11.11}
两个绝热容器各有 \(\nu\,\mathrm{mol}\)
同种理想气体，起初温度相同，
压力分别为 \(p_1,p_2\)。
接通容器并充分平衡，
证明系统熵变
\[
 \Delta S=\nu R
 \ln\frac{(p_1+p_2)^2}{4p_1p_2}
\]
且 \(\Delta S\ge0\)。
\textit{容器保持刚性，气体的化学组成相同。}

\paragraph{习题 11.12}
一块质量 \(m\)、定压比热 \(c_p\)、
初温 \(T_i\) 的金属，
放入外界绝热、压强不变的液体中。
液体质量 \(m'\)、比热 \(c'_p\)、初温 \(T'_i\)。
求最终共同温度，证明最后达到热平衡时
二者合系统的熵为最大值，
并给出整个过程的熵变。

\paragraph{习题 11.13}
根据图11.16的\emph{氮气温熵图}
（纵轴 \(T/\mathrm K\)，横轴
\(S_m/[\mathrm{J/(mol\,K)}]\)，其中有等压线
与等焓线）作图估计：
（1）氮气从 \(80\,\mathrm K,50\,\mathrm{MPa}\)
节流至 \(20\,\mathrm{MPa}\) 的末温；
（2）从 \(70\,\mathrm K,2.0\,\mathrm{MPa}\)
节流至 \(0.1\,\mathrm{MPa}\) 的末温；
（3）从 \(76\,\mathrm K,5.0\,\mathrm{MPa}\)
作\emph{可逆绝热}膨胀至
\(0.1\,\mathrm{MPa}\) 的末温。
\textit{节流应按等焓 \(H_m=\mathrm{const}\)，
可逆绝热应按等熵 \(S_m=\mathrm{const}\)；
图读温度需保留画图误差与配套答案差异。}
